\documentclass[12pt]{article}
\usepackage[utf8]{inputenc}
\usepackage[T1]{fontenc}
\usepackage{hyperref}
\title{Kaur Language Resources for K-12 Curriculum Development}
\author{Compiled by Claude Sonnet 5 (claude-sonnet-5) for Andr\'as Kornai}
\date{}
\begin{document}
\maketitle

\section{Introduction}
Kaur is a Malayic language of Kaur Regency, Bengkulu Province, on the southwest coast of Sumatra. Wikipedia cites about 40,000 speakers (2000 figure). Dialects include South Kaur (around Bintuhan, influenced by Lampung) and North Kaur (influenced by Besemah). Speaker figures are dated, so current numbers are unclear.

\section{Primary Educational Resources}
I found no Kaur textbooks, primers or school programmes. Schooling is in Indonesian.

\section{Linguistic and Reference Materials}
\begin{itemize}
\item Kaur language (Wikipedia): \url{https://en.wikipedia.org/wiki/Kaur_language}
\item Wikidata entry for Kaur: \url{https://www.wikidata.org/wiki/Q3182322}
\item Documentation of multilingual practices in the Nasal speech community, Bengkulu: \url{https://researchdata.edu.au/documentation-multilingual-linguistic-speech-community/1467645}
\item Learn 100 everyday Bengkulu words, Kaur dialect (Good News from Indonesia, Indonesian): \url{https://www.goodnewsfromindonesia.id/2025/04/22/belajar-100-kosakata-bahasa-bengkulu-sehari-hari-dialek-kaur}
\end{itemize}

\section{Religious and Community Resources}
\begin{itemize}
\item Introduction to the Kaur people of Bengkulu (detikSumbagsel, Indonesian): \url{https://www.detik.com/sumbagsel/budaya/d-7167055/mengenal-suku-kaur-di-bengkulu-asal-usul-hingga-mata-pencaharian}
\item Asia Harvest people-group profile: \url{https://people-groups.asiaharvest.org/Indonesia/Kaur.pdf}
\item Praying for Indonesia entry: \url{https://www.prayingforindonesia.com/?p=954}
\end{itemize}

\section{Digital and Audio Resources}
No dedicated audio, video or software resources were found beyond the documentation collection listed above.

\section{Curriculum Development Context}
Kaur is vigorous in its community but not known to be taught in schools. Materials would have to be built from scratch, starting from the documentation and the vocabulary list above. Local-content (muatan lokal) provisions in Indonesian schools might be a route, but I did not search for them.
\end{document}